\clearpage
\VerdierPhysicalPageAnchor{0178}%
\begingroup\small\setlength{\parskip}{0pt}\setlength{\emergencystretch}{1.5em}
\setlength{\abovedisplayskip}{4pt plus 2pt minus 2pt}%
\setlength{\belowdisplayskip}{4pt plus 2pt minus 2pt}%

\VerdierSemanticAnchor{ch3d-1-2-3-a-cont}
\noindent\textit{sont strictement pleinement fidèles. Ils sont injectifs sur les
objets et réalisent donc les catégories $K^{,*}(\mathcal A)/\operatorname{Ac}(\mathcal A)$
comme sous-catégories triangulées strictement pleines de $D(\mathcal A)$. Un
objet $X$ de $D(\mathcal A)$ est un objet de
$K^{,b}(\mathcal A)/\operatorname{Ac}(\mathcal A)$ (resp.
$K^{,+}(\mathcal A)/\operatorname{Ac}(\mathcal A)$, resp.
$K^{,-}(\mathcal A)/\operatorname{Ac}(\mathcal A)$) si et seulement si
$H^i(X)=0$, pour $|i|$ grand (resp. $i$ petit, resp. $i$ grand).}

\VerdierSemanticAnchor{ch3d-1-2-3-b-diagrams}
\noindent\textit{b) Les foncteurs naturels qui figurent en traits pleins dans les
diagrammes :}
\begin{center}
\resizebox{0.80\textwidth}{!}{%
$\begin{tikzcd}[ampersand replacement=\&,column sep=4.2em,row sep=2.7em]
K^{+}(\mathcal A)/\operatorname{Ac}^{+}(\mathcal A)
  \arrow[r] \& K^{,+}(\mathcal A)/\operatorname{Ac}(\mathcal A)
  \arrow[dr,dashed] \& \\
\& \& D(\mathcal A) \\
K^{-}(\mathcal A)/\operatorname{Ac}^{-}(\mathcal A)
  \arrow[r] \& K^{,-}(\mathcal A)/\operatorname{Ac}(\mathcal A)
  \arrow[ur,dashed] \&
\end{tikzcd}$%
}

\medskip
\resizebox{0.88\textwidth}{!}{%
$\begin{tikzcd}[ampersand replacement=\&,column sep=3.2em,row sep=2.1em]
\& K^{+,b}(\mathcal A)/\operatorname{Ac}^{+}(\mathcal A) \arrow[dr] \& \& \\
K^{b}(\mathcal A)/\operatorname{Ac}^{b}(\mathcal A)
  \arrow[ur] \arrow[dr]
  \& \& K^{,b}(\mathcal A)/\operatorname{Ac}(\mathcal A)
  \arrow[r,dashed] \& D(\mathcal A) \\
\& K^{-,b}(\mathcal A)/\operatorname{Ac}^{-}(\mathcal A) \arrow[ur] \& \&
\end{tikzcd}$%
}
\end{center}
\noindent\textit{sont des équivalences de catégories. Ils sont injectifs sur les
ensembles d'objets.}

\VerdierSemanticAnchor{ch3d-1-2-3-c-diagram}
\noindent\textit{c) Les foncteurs naturels qui figurent en traits pleins dans le
diagramme :}
\begin{center}
\resizebox{0.82\textwidth}{!}{%
$\begin{tikzcd}[ampersand replacement=\&,column sep=4.0em,row sep=2.7em]
K^{+,b}(\mathcal A)/\operatorname{Ac}^{+}(\mathcal A)
  \arrow[r] \& K^{+}(\mathcal A)/\operatorname{Ac}^{+}(\mathcal A)
  \arrow[dr,dashed] \& \\
\& \& D(\mathcal A) \\
K^{-,b}(\mathcal A)/\operatorname{Ac}^{-}(\mathcal A)
  \arrow[r] \& K^{-}(\mathcal A)/\operatorname{Ac}^{-}(\mathcal A)
  \arrow[ur,dashed] \&
\end{tikzcd}$%
}
\end{center}
\noindent\textit{sont strictement pleinement fidèles et injectifs sur les
ensembles d'objets.}

\VerdierSemanticAnchor{ch3d-1-2-3-proof}
Résulte immédiatement de (1.2.1) et de (chap. II, 2.3.5).
\endgroup
